\documentclass{../base/canon}

%% canon_variables.tex — Valores por defecto de metadatos institucionales
%%
%% Incluir ANTES del \begin{document} para sobreescribir los defaults de la clase.
%% El template puede sobreescribir cualquier valor después de este \input.
%%
%% Uso en template:
%%   %% canon_variables.tex — Valores por defecto de metadatos institucionales
%%
%% Incluir ANTES del \begin{document} para sobreescribir los defaults de la clase.
%% El template puede sobreescribir cualquier valor después de este \input.
%%
%% Uso en template:
%%   %% canon_variables.tex — Valores por defecto de metadatos institucionales
%%
%% Incluir ANTES del \begin{document} para sobreescribir los defaults de la clase.
%% El template puede sobreescribir cualquier valor después de este \input.
%%
%% Uso en template:
%%   \input{../base/canon_variables}
%%   \SetDocTitle{Mi Informe Específico}   % sobreescribe el default

\SetDocLogo{../assets/logo.pdf}
\SetDocUnit{Tecnología \& Sistemas}
\SetContactUnit{Mesa de soporte}
\SetContactMembers{%
  Equipo CoreX & Soporte & soporte@empresa.com%
}
\SetContactNote{Para soporte técnico o consultas sobre este documento.}

%%   \SetDocTitle{Mi Informe Específico}   % sobreescribe el default

\SetDocLogo{../assets/logo.pdf}
\SetDocUnit{Tecnología \& Sistemas}
\SetContactUnit{Mesa de soporte}
\SetContactMembers{%
  Equipo CoreX & Soporte & soporte@empresa.com%
}
\SetContactNote{Para soporte técnico o consultas sobre este documento.}

%%   \SetDocTitle{Mi Informe Específico}   % sobreescribe el default

\SetDocLogo{../assets/logo.pdf}
\SetDocUnit{Tecnología \& Sistemas}
\SetContactUnit{Mesa de soporte}
\SetContactMembers{%
  Equipo CoreX & Soporte & soporte@empresa.com%
}
\SetContactNote{Para soporte técnico o consultas sobre este documento.}

\SetDocType{INFORME EJECUTIVO}
\SetDocTitle{Rendimiento Operativo Postcosecha — Semana 17-2026}
\SetDocCode{INF-OPS-017}
\SetDocArea{Operaciones / Postcosecha}
\SetDocAuthor{Jefe de Operaciones}
\SetDocDate{2026-04-22}
\SetContactUnit{Jefatura de Operaciones}
\SetContactMembers{%
  Carlos Mendoza   & Jefe de Operaciones  & c.mendoza@empresa.com   \\
  Patricia Suárez  & Calidad Postcosecha  & p.suarez@empresa.com%
}
\SetContactNote{Para consultas sobre este informe contactar antes de las 16:00.}

\begin{document}

\section{Resumen ejecutivo}

\begin{ParrafoEjecutivo}
La semana 17-2026 cerró con \textbf{842\,000 tallos procesados} y un rendimiento
general de clasificación del \textbf{96.3\,\%}, superando el objetivo del 95.0\,\%.
La merma total se mantuvo en 2.1\,\%, dentro del rango aceptable.
No se registraron paros de línea no programados ni incidentes de seguridad industrial.
\end{ParrafoEjecutivo}

\section{Indicadores clave}

\begin{center}
\begin{tabular}{ccc}
  \FichaKPI{Tallos procesados}{842\,000}{semana 17-2026} &
  \FichaKPI{Rendimiento general}{96.3\,\%}{meta: 95.0\,\%} &
  \FichaKPI{Merma total}{2.1\,\%}{límite: 3.0\,\%} \\[1.4em]
  \FichaKPI{Lotes completados}{15 / 15}{100\,\% programados} &
  \FichaKPI{Ciclo promedio}{4.2 min}{meta: ≤ 4.5 min} &
  \FichaKPI{Paros no prog.}{0}{sin incidentes} \\
\end{tabular}
\end{center}

\section{Resultados por área}

\begin{table}[H]
  \centering
  \caption{Rendimiento por área de proceso — Semana 17-2026}
  \begin{tabular}{l r r r r}
    \toprule
    \textbf{Área} & \textbf{Ingreso} & \textbf{Salida} & \textbf{Merma} & \textbf{Rend.\,\%} \\
    \midrule
    Apertura       & 280\,400 & 271\,300 & 9\,100 & 96.8 \\
    Clasificación  & 271\,300 & 261\,200 & 10\,100 & 96.3 \\
    Empaque        & 261\,200 & 257\,900 & 3\,300 & 98.7 \\
    \midrule
    \textbf{Total} & \textbf{812\,900} & \textbf{790\,400} & \textbf{22\,500} & \textbf{97.2} \\
    \bottomrule
  \end{tabular}
\end{table}

\section{Hallazgos}

\begin{itemize}
  \item El rendimiento de empaque alcanzó 98.7\,\%, el mejor valor del ciclo 2026.
  \item Las variedades Freedom y Vendela concentraron el 68\,\% de la merma total
        por causa «botón abierto».
  \item El turno tarde del día 15-abr presentó un incremento puntual de merma del
        18\,\% respecto al promedio de la semana; se identificó y corrigió en turno.
\end{itemize}

\begin{ObservationBox}[Observación de calidad]
  Dos lotes (LOT-2026-1408, LOT-2026-1411) registraron merma por «botrytis» superior
  al 3.0\,\%. Se notificó al área de campo para revisión de condiciones de cosecha
  en las zonas de origen correspondientes.
\end{ObservationBox}

\section{Recomendaciones}

\begin{enumerate}
  \item Implementar revisión adicional de madurez en lotes de variedad Freedom antes
        de ingreso a clasificación — \textit{Responsable: Supervisión Turno Mañana,
        a partir de la semana 18}.
  \item Evaluar ajuste del umbral de merma por «botón abierto» para variedades de
        tallo corto — \textit{Responsable: Jefatura de Calidad, plazo: 2026-04-30}.
  \item Revisar condiciones de cadena de frío en origen de los lotes con botrytis —
        \textit{Responsable: Coordinación de Campo, plazo: 2026-04-25}.
\end{enumerate}

\SignatureBlock{Carlos Mendoza}{Jefe de Operaciones Postcosecha}

\ContactSignature

\end{document}
